\section{Study Design and Evidence Scope}
\label{sec:methods}
\subsection{Frozen systems and complete endpoints}
Each study has its own protocol, tasks and comparator, stated beside its results and mapped in the historical appendix and the separate intervention appendix. The target/draft pair is Qwen3-8B/0.6B Q8\_0, except the 358H adapter's symmetric INT4 OpenVINO draft. Final comparisons process one request at a time. Historical settings are in Appendix K and RTX V4 in Appendix Q; hardware, backend, OS and quantization jointly vary, preventing hardware-only causal comparisons. RTX's original direct-ID OS build is 26200; follow-ups use 26300. A driver version alone does not identify driver bytes.

``Complete'' retains each protocol's generated-ID and endpoint rules: prompt IDs and uncommitted proposals are excluded; an earlier terminal EOG is inclusive, otherwise the 128-ID cap applies. Agreement compares all IDs within that endpoint and the endpoint itself. B570 raw native captures additionally retain 1--2 committed post-cap IDs, excluded from agreement but exported for audit. Actual stop markers are checked. Valid results are selected from each result manifest's complete attempt; failed or interrupted attempts are retained separately and are never concatenated.

The B570 intervention uses a previously exposed 24-task diagnostic/control cohort. RTX uses the old 19-task cohort and a separately frozen 24-task MBPP cohort; its candidate, exposure exclusions and static task order precede new-task execution. Frozen prediction arrays precede native runs. RTX additionally records remote readback before native; B570 uses local seals. Neither provides independently trusted absolute experimental time. Prospective prediction is not equivalent to formal independent preregistration, and the candidate's earlier exposure is recorded. No tasks are added or dropped after observing these results. Three repeats assess within-task stability; finite exact counts receive no pseudo-independent binomial intervals.

The offline intervention review restores the hash-bound evidence archive and independently parses selected result/raw/trace records. It checks prediction/native and native/reference arrays and endpoints, repeat stability, paired transitions, handoff witnesses, shams, graph occupancy and effective settings. Opaque sampler state cannot be serialized: clone behavior and checked surrounding state do not establish byte equality of opaque sampler bytes. Omitted weights/runtime binaries remain identified but are not independently rehashed or executed by this review. The author evidence packet records exact review coverage and source/hash bindings.

Historical M4 includes a representative 24-prompt Spec-Bench sample \citep{specbench} and a separate preregistered 30-task HumanEval study \citep{humaneval}; they are background, not replications of the later runtime study. The latter retains its own calibration, qualification and held-out splits in \ifruntime Section~\ref{sec:runtime}\else the full manuscript\fi. All studies keep their original stop rules and comparators.

\subsection{Cost estimators and conditional uncertainty}
Later RTX cost studies take each condition's median of three times per prompt, then reference/candidate ratios and their prompt GM. V4 uses 10,000 paired-row bootstrap draws, seed 2026100204, frozen task order and linear .025/.975 percentiles. This resamples the empirical distribution of the selected prompt rows, conditional on the exposure-audited cohort, observed condition medians and frozen configuration. It does not account for exposure selection, predict a HumanEval population failure rate or cover a new thermal/session distribution. Repetitions are not independent bootstrap units. Cost interpretation requires every reference-ID pair and stable reference repeats; blocked timing remains diagnostic.

The 358H prompt estimator takes the median of three paired ratios before the prompt GM; total-time accounting is separate. M4 runtime uses condition medians and its own frozen random streams. Historical H estimators remain in Appendix B. All intervals omit full session and recalibration uncertainty.

Timing excludes loading and initial prefill. The RTX V4 timer begins after synchronized target/draft prefill and before per-request sampler/state work, ending at slot release. The M4 and 358H boundaries differ and are stated with their results. M4 package energy excludes separate DRAM, B570 uses adapter-board energy, and RTX mixed domains do not form a common system boundary. There is no end-to-end 358H Phase-3 or runtime energy result.

\subsection{Offline reproducibility}
The anonymous index recomputes retained arrays, timer intervals and summaries; the author packet supplies source/hash bindings and raw-review coverage. Appendix O maps historical entries and Q the V4 capabilities. Verification runs no models and does not authenticate omitted assets or absolute process time.

\paragraph{Existing-record analysis.}
Score/margin descriptors, trajectory extents, fixed-mode opportunity and conditional cost bounds are calculated after aggregate outcomes are known. Extent/budget analyses retain original rows and comparators without adding samples or intervals. V4 latency summaries describe 19 prompt medians per condition without changing the paired GM/CI. Appendix P gives the derivations.
